\section{From local cohomology to divisorial valuations}
\label{sec:local-valuations}

We prove the local assertion used in Theorem~\ref{thm:branch}.  Throughout
this section the ground field is $\mathbb C$.  A fibre is its scheme fibre
unless explicitly given the reduced structure.  In particular, line
bundles and coherent cohomology on a resolution fibre are taken on that
scheme, including its possible nilpotents.

\begin{proposition}\label{prop:local-valuations}
Let $Y$ be a normal projective $\mathbb Q$-factorial fourfold, and let
$(Y,\Theta)$ be an effective terminal pair with rational boundary.  Let
$T$ be an irreducible curve in $\operatorname{Sing}Y$ contained in a
positive-coefficient component of $\Theta$.  Put
\[
 j:Y_{\mathrm{reg}}\hookrightarrow Y,
 \qquad
 \mathcal V_T=(R^2j_*\mathbb Q)|_{T^\circ},
 \qquad
 b_T=\dim_{\mathbb Q}H_c^2(T^\circ,\mathcal V_T),
\]
where $T^\circ$ is a sufficiently small smooth dense open on which the
indicated sheaf is a local system.  There are at least $b_T$ distinct
normalized divisorial valuations $v$ of $\mathbb C(Y)$ such that
\[
 c_Y(v)=T,\qquad a(v;Y,0)=2,\qquad a(v;Y,\Theta)<2.
\]
The valuations obtained for different curves $T$ are distinct.
\end{proposition}

The argument has two descent steps.  Norms first produce line bundles
over the function field of $T$ itself.  Approximation then extends those
bundles to a pointed \emph{ordinary} \'etale
neighbourhood, with the same residue field at the chosen point.  This
last condition is needed to keep the resulting valuations distinct.
On a small $\mathbb Q$-factorial model of this neighbourhood, the extended
bundles give independent degree vectors on the irreducible curve components
of the fibre over the chosen point.  These components are counted over the
function field of $T$ (Lemma~\ref{local:degree-independence}).
Blowing up their surface closures at their generic points gives
log-discrepancy-two valuations, and Lemma~\ref{local:restriction-injective}
makes their restrictions to $\mathbb C(Y)$ distinct.

\subsection{Resolution fibres and their monodromy}

Fix a projective resolution
\[
 \pi:\widetilde Y\longrightarrow Y
\]
which is an isomorphism over $Y_{\mathrm{reg}}$.  Since terminal
singularities are smooth in codimension two, $\dim\operatorname{Sing}Y
\leq1$.  After shrinking $T^\circ$, the fibres of $\pi$ over $T^\circ$
have dimension at most two.  Indeed, a component of $\pi^{-1}(T)$
dominating $T$ has dimension at most three, because it is a proper
closed subset of the fourfold $\widetilde Y$.  Upper semicontinuity then
removes the finitely many possible dimension jumps.  We also shrink so
that
\[
 \mathcal F=\widetilde Y\times_Y T^\circ\longrightarrow T^\circ
\]
is flat and projective and
$\mathcal W=(R^2\pi_*\mathbb Q)|_{T^\circ}$ is a local system.

\begin{lemma}\label{local:split-punctured}
Restriction to the smooth open defines a split surjection of rational
local systems
\[
 r:\mathcal W\longrightarrow\mathcal V_T.
\]
The splitting is available on $T^\circ$ itself, before passing to any
cover.
\end{lemma}

\begin{proof}
Use the perverse normalization for $\mathrm{IC}_Y$, so that its restriction to
$Y_{\mathrm{reg}}$ is $\mathbb Q[4]$.  The decomposition theorem for the
projective morphism $\pi$ provides an inclusion of the full-support
summand
\[
 s:\mathrm{IC}_Y\longrightarrow R\pi_*\mathbb Q[4]
\]
whose restriction to $Y_{\mathrm{reg}}$ is the identity; see
\cite[Theorem~5.3.1]{Saito1988} and
\cite[\S\S1.7--1.13]{Saito1989}.
Restriction upstairs gives a natural morphism
\[
 R\pi_*\mathbb Q[4]\longrightarrow Rj_*\mathbb Q[4].
\]
Its composite with $s$ is the canonical comparison
$\mathrm{IC}_Y\to Rj_*\mathbb Q[4]$.  In fact, adjunction identifies morphisms
from $\mathrm{IC}_Y$ to $Rj_*\mathbb Q[4]$ with morphisms from
$j^*\mathrm{IC}_Y=\mathbb Q[4]$ to itself, and the indicated composite restricts
to the identity.

At the generic curve stratum, which has complex codimension three,
the middle-extension construction retains the ordinary unshifted
cohomology through degree two.  Consequently the comparison induces
\[
 \mathcal H^{-2}(\mathrm{IC}_Y)|_{T^\circ}
   \xrightarrow{\ \sim\ }(R^2j_*\mathbb Q)|_{T^\circ}.
\]
This is the middle-perversity truncation description of intersection
cohomology \cite[\S2]{GM1983}; subsequent extension across the
zero-dimensional strata does not change this restriction.  Taking
degree $-2$ in the two morphisms above proves the assertion: the map
induced by $s$, composed with the inverse of this comparison
isomorphism, is a right inverse to $r$.
\end{proof}

\begin{lemma}\label{local:scheme-fibre-vanishing}
Let $y\in Y$ be a point for which the scheme fibre
$F_y=\widetilde Y\times_Y\operatorname{Spec}\kappa(y)$ has dimension at
most two.  Then
\[
 H^2(F_y,\mathcal O_{F_y})=0.
\]
If $y$ is a closed complex point, first Chern classes of algebraic line
bundles span $H^2(F_y,\mathbb Q)$.
\end{lemma}

\begin{proof}
The variety $Y$ has rational singularities, since it is klt
\cite[Theorem~5.22]{KM98}.  In particular,
$R^2\pi_*\mathcal O_{\widetilde Y}=0$.  Write
$A=\mathcal O_{Y,y}$, let $\mathfrak m$ be its maximal ideal, and put
\[
 F_n=\widetilde Y\times_Y
       \operatorname{Spec}(A/\mathfrak m^{n+1})\qquad(n\geq0).
\]
Every $F_n$ has the same underlying topological space as $F_0=F_y$.
In the exact sequence
\[
 0\longrightarrow J_n\longrightarrow\mathcal O_{F_{n+1}}
   \longrightarrow\mathcal O_{F_n}\longrightarrow0,
\]
the coherent kernel has $H^3(J_n)=0$, by the dimension bound.
Therefore
$H^2(F_{n+1},\mathcal O_{F_{n+1}})\to
H^2(F_n,\mathcal O_{F_n})$ is surjective.  Proper formal functions
\cite[Tag~02OD]{Stacks} gives
\[
 \varprojlim_n H^2(F_n,\mathcal O_{F_n})
   =(R^2\pi_*\mathcal O_{\widetilde Y})_y^{\wedge}=0.
\]
For an inverse sequence of surjective maps the limit maps
surjectively to each term, by successive lifting.  Hence every term,
in particular $H^2(F_y,\mathcal O_{F_y})$, is zero.  This argument
applies also when $y$ is the generic point of $T$; no flatness at $y$
is required.

For closed $y$, the analytic exponential sequence on the projective
complex scheme $F_y$ gives
\[
 \operatorname{Pic}(F_y^{\mathrm{an}})
  \xrightarrow{\ c_1\ } H^2(F_y^{\mathrm{an}},\mathbb Z)
  \longrightarrow H^2(F_y^{\mathrm{an}},\mathcal O_{F_y^{\mathrm{an}}}).
\]
The exponential sequence remains exact in the presence of nilpotents:
the exponential and logarithm are inverse on a nilpotent ideal, and
the usual sequence is exact on the reduction.  The scheme form of
GAGA \cite[Corollary~4.3 and Theorem~4.4]{RaynaudSGA1XII},
extending Serre's projective comparison \cite{SerreGAGA},
identifies the last group with the group already proved to vanish
and identifies analytic line bundles with algebraic ones, including
on this possibly nonreduced proper scheme. Thus $c_1$ is surjective
integrally, and tensoring with $\mathbb Q$ proves the assertion.
\end{proof}

\begin{lemma}\label{local:generic-lines}
The local systems $\mathcal W$ and $\mathcal V_T$ have finite
monodromy.  Put $K_T=\kappa(\eta_T)$ and
$F_{\eta_T}=\widetilde Y\times_Y\operatorname{Spec}K_T$.  There are
$b_T$ line bundles
\[
 L_1,\ldots,L_{b_T}\in\operatorname{Pic}(F_{\eta_T})
\]
whose spreads over a sufficiently small $T^\circ$ have linearly
independent images under $r$ in every fibre of $\mathcal V_T$.
\end{lemma}

\begin{proof}
We give the parameter argument because the bundles are needed over
$K_T$, rather than only over its algebraic closure.  Fix a relatively
ample line bundle $\mathcal O_{\mathcal F}(1)$.  Every line bundle on
a closed fibre is a quotient of some
$\mathcal O_{F_t}(-m)^{\oplus N}$.  Let $m$, $N$, and the Hilbert
polynomial vary. Grothendieck's Quot-scheme theorem
\cite[Theorems~3.1--3.2]{GrothendieckHilbert1961} gives the
corresponding countable collection of schemes of finite type over $T^\circ$.
On each, take the open locus where the universal quotient is
invertible on the whole pulled-back fibre family.  These loci still
parameterize every line bundle on every closed fibre.  To justify
openness, the family and the universal quotient are flat over the
Quot scheme; a fibrewise locally free quotient of rank one is locally
free near that fibre by the local criterion for flatness.  Properness
then removes the image of the complementary closed locus.

There are only countably many irreducible components of these
parameter spaces.  Choose a very general $t\in T^\circ(\mathbb C)$
outside the closures of the images of all components which do not
dominate $T^\circ$.  Lemma~\ref{local:scheme-fibre-vanishing} supplies
finitely many line bundles on $F_t$ whose classes
$c_1,\ldots,c_N$ span $\mathcal W_t$.  Each is represented by a point
of a dominating irreducible parameter component.  The first Chern
class of the universal line bundle gives a section of the pullback
of $\mathcal W$ to that component: this follows by restricting its
class on the total family to the fibres and applying proper base
change.  The analytic space of this irreducible complex parameter
component is path connected.

Choose a closed point of the generic fibre of this parameter
component.  Its residue field is finite over $K_T$.  Spreading the
point, normalizing, and shrinking the base yields a connected finite
\'etale cover of a dense open of $T^\circ$, together
with a line bundle on the pulled-back family. Choose a path in the
parameter component from the original quotient point to a point
on this cover. Flatness of the universal Chern class identifies
the transported $c_i$ with the class at that cover point. The
finite-index subgroup fixing the chosen sheet fixes this class;
conjugating by the projected path gives a finite-index subgroup
fixing $c_i$ at the original basepoint. Shrinking a smooth complex curve does
not diminish the monodromy image of a local system originally
defined on it: the induced map on fundamental groups is surjective.
It follows that each $c_i$ has a finite orbit under
$\Gamma=\pi_1(T^\circ,t)$.  Intersecting the stabilizers of this
finite spanning set gives a subgroup of finite index acting trivially
on $\mathcal W_t$.  The image $G$ of $\Gamma$ in
$\operatorname{GL}(\mathcal W_t)$ is therefore finite.
Lemma~\ref{local:split-punctured} gives the same conclusion for
$\mathcal V_T$.
Take a common dense open where all these covers are finite
\'etale, move the basepoint into it, and transport
the $c_i$ along one common path to the new basepoint, still denoted
$t$.  They still span $\mathcal W_t$.  Each cover class is a
monodromy translate of the corresponding transported $c_i$.
Changing its sheet realizes every element of that orbit, so choose
a sheet giving $c_i$ itself.

We next take norms of actual line bundles, not merely invariant
points of a Picard functor.  For the connected cover associated with
$c_i$, let $H_i\subset\Gamma$ be the subgroup fixing its chosen
sheet, and let $d_i=[\Gamma:H_i]$.  The cover carries a line bundle
$\mathcal L_i$ whose class at that sheet is $c_i$ and is fixed by $H_i$.  Its norm
along the finite locally free map of fibre families is a line bundle
on the original family; norms commute with base change
\cite[Tags~0BCY and~0BD2]{Stacks}.  At $t$ its first Chern class is
\begin{equation}\label{eq:local-norm-average}
 c_1(\operatorname{Nm}\mathcal L_i)
   =\sum_{\gamma\in\Gamma/H_i}\gamma c_i
   =d_i\operatorname{Av}_G(c_i),
 \qquad
 \operatorname{Av}_G(c)=\frac1{|G|}\sum_{g\in G}gc.
\end{equation}
For completeness, $H_i$ need not contain the kernel of
$\Gamma\to G$.  Nevertheless its image fixes $c_i$, and each element
of the finite $G$-orbit of $c_i$ occurs equally often in the coset
sum.  Both expressions in \eqref{eq:local-norm-average} are
$d_i$ times the orbit average, which proves the equality with its
stated factor.

The averages of a spanning set span $\mathcal W_t^G$.  Moreover,
the surjection $r$ induces a surjection on invariants: averaging any
lift of an invariant vector gives an invariant lift.  Poincar\'e
duality on the connected oriented real surface $T^\circ$ identifies
the dimension of $H_c^2(T^\circ,\mathcal V_T)$ with the dimension of
the coinvariants of its monodromy representation.  For a finite
group over $\mathbb Q$, invariants and coinvariants have the same
dimension, since averaging identifies the latter with the former.
Thus
\[
 b_T=\dim(\mathcal V_{T,t})^G.
\]
The norm bundles have generic fibres defined over $K_T$ itself, and
their images span these invariants.  Select $b_T$ of them whose
images form a basis and rename their generic fibres $L_1,\ldots,L_{b_T}$.
On a common smaller open, their images are independent flat
sections, proving the claim.  Their invariant dimension, and hence
$b_T$, is unchanged by further shrinking.
\end{proof}

\subsection{Extending line bundles without changing the residue field}

The independent punctured classes are now represented by line bundles
on the generic scheme fibre over $K_T$.  We next extend these bundles
in transverse directions while keeping that residue field unchanged.

\begin{lemma}\label{local:etale-lines}
Let $L_1,\ldots,L_b$ be finitely many line bundles on $F_{\eta_T}$.
There exist a normal integral quasi-projective variety $Y_e$, an
\'etale morphism $q:Y_e\to Y$, and a point
$\eta'\in Y_e$ above $\eta_T$ such that
\[
 K_T=\kappa(\eta_T)\xrightarrow{\ \sim\ }\kappa(\eta')
\]
is the residue-field map, and every $L_i$ extends to a line bundle
$\widetilde L_i$ on
$\widetilde Y_e=\widetilde Y\times_Y Y_e$.  The extensions restrict
to the specified $L_i$ on the same scheme fibre.  The variety
$\widetilde Y_e$ is smooth and integral, and
$\widetilde Y_e\to Y_e$ is projective and birational.
\end{lemma}

\begin{proof}
Write $A=\mathcal O_{Y,\eta_T}$, $\mathfrak m=\mathfrak m_A$, and
$Z=\widetilde Y\times_Y\operatorname{Spec}A$.  The ring $A$ is an
excellent local ring essentially of finite type over $\mathbb C$,
with residue field $K_T$.  Let $\widehat A$ be its completion, and
write $F_n=Z\times_A A/\mathfrak m^{n+1}$ as in the proof of
Lemma~\ref{local:scheme-fibre-vanishing}.

For $n\geq0$, the square-zero ideal for $F_n\subset F_{n+1}$ is
\[
 J_n=\mathfrak m^{n+1}\mathcal O_Z/
                       \mathfrak m^{n+2}\mathcal O_Z.
\]
Multiplication gives a surjection
\[
 (\mathfrak m^{n+1}/\mathfrak m^{n+2})
       \otimes_{K_T}\mathcal O_{F_0}\twoheadrightarrow J_n.
\]
Its source is a finite direct sum of $\mathcal O_{F_0}$.  The kernel
is coherent and has vanishing third cohomology, since $\dim F_0\leq2$.
Lemma~\ref{local:scheme-fibre-vanishing} therefore implies
$H^2(J_n)=0$.  The exact sequence of units for a square-zero
thickening, with $1+J_n\simeq J_n$ additively, shows that the
obstruction to lifting a line bundle from $F_n$ to $F_{n+1}$ lies in
this zero group \cite[Tag~0C6R]{Stacks}.  Each $L_i$ consequently has
compatible lifts to every $F_n$.

Grothendieck's existence theorem \cite{EGAIII1}, in the projective
form \cite[Tag~0885]{Stacks}, for the scheme $Z_{\widehat A}$
algebraizes these compatible systems of coherent sheaves
and their morphisms. It produces invertible sheaves: algebraize
the systems of inverse line bundles at the same time, and use full
faithfulness to algebraize their tensor-product isomorphisms with
$\mathcal O$.  Denote the resulting line bundles on
$Z_{\widehat A}$ by $\widehat L_i$.  Their restrictions to $F_0$
are the specified $L_i$, with the chosen identifications.

\begingroup\emergencystretch=1em
The finite-presentation limit property for invertible sheaves and
their isomorphisms \cite[Tags~0B8W and~01ZR]{Stacks} descends all
$\widehat L_i$ to line bundles on $Z_R$ for a single finitely
generated $A$-algebra $R$ equipped with a map
$\phi:R\to\widehat A$.  One may include their inverses and the
evaluation isomorphisms in these data.  Since $A$ is noetherian,
$R$ has a finite presentation.  Applying Artin approximation to its
finite system of equations gives a map
\[
 \psi:R\longrightarrow A^h
\]
whose reduction in $K_T$ agrees with that of $\phi$, where $A^h$ is
the ordinary henselization of $A$, whose completion is $\widehat A$.
Only approximation modulo
$\mathfrak m$ is used.  This is the approximation theorem for an
essentially finite-type local ring over a field
\cite[Theorem~1.10]{Artin1969}; the pointed \'etale
formulation, including equality of residue fields, is
\cite[Tag~07QZ]{Stacks}.
\par\endgroup

Pulling the descended line bundles back by $\psi$ gives line bundles
on $Z_{A^h}$ with the required special fibres.  More explicitly, the
two maps $R\to K_T$ are identical, so their restrictions are the
same pullbacks of the line bundles on $Z_R$ to $Z_{K_T}=F_0$.
Thus the argument preserves bundles on the entire scheme fibre;
it does not merely preserve their numerical classes.

The ordinary henselization is the filtered colimit of pointed
\'etale neighbourhoods with unchanged residue
field \cite[Tags~0A02 and~0A03]{Stacks}.  The same
finite-presentation property descends our finite list of bundles to
one such neighbourhood.  Spreading it from $\operatorname{Spec}A$
to an affine open of $Y$ gives $q:Y_e\to Y$ and $\eta'$ as claimed,
after shrinking around $\eta'$.  Since $Y$ is normal, its
\'etale neighbourhood is normal.  Its irreducible
components are open and closed; take the one containing $\eta'$.
It is integral and dominant over $Y$, since an
\'etale map is open.  We may take $Y_e$
quasi-projective by the affine construction.

Finally $\widetilde Y_e$ is smooth, since it is
\'etale over the smooth variety $\widetilde Y$.
Its irreducible components are disjoint and open.  Every component
has nonempty open image in $\widetilde Y$, and hence meets
$\pi^{-1}(Y_{\mathrm{reg}})$.  Over this dense open the base change
is the nonempty integral open $q^{-1}(Y_{\mathrm{reg}})$ of $Y_e$.
Every component meeting that integral open forces there to be just
one component.  Hence $\widetilde Y_e$ is integral.  Its map to
$Y_e$ is projective and is an isomorphism over
$q^{-1}(Y_{\mathrm{reg}})$, so it is birational.
\end{proof}

Apply the lemma to the bundles in Lemma~\ref{local:generic-lines},
put $b=b_T$, and let $T'\subset Y_e$ be the closure of $\eta'$.
Write $\pi_e:\widetilde Y_e\to Y_e$ for the base-changed resolution and
$j_e:(Y_e)_{\mathrm{reg}}\hookrightarrow Y_e$ for the smooth-locus inclusion.
The induced map
$T'\to T$ is an isomorphism on suitable dense opens: it is
generically an isomorphism by the residue-field identity.  On that
open the restrictions of $\widetilde L_i$ agree with the spreads used
in Lemma~\ref{local:generic-lines}.  Indeed, their isomorphisms on the
generic scheme fibre spread after shrinking, by finite presentation.
At a closed complex point an \'etale map is a
local analytic isomorphism.  The local systems of punctured
cohomology and the maps from resolution-fibre cohomology are
therefore identified near these points.  It follows that the images
of $c_1(\widetilde L_1),\ldots,c_1(\widetilde L_b)$ in punctured
degree-two cohomology along a dense open of $T'$ remain independent.

\subsection{Degrees on a small model}

Although $Y$ is $\mathbb Q$-factorial, its \'etale neighbourhood
$Y_e$ need not be.  We pass to a small $\mathbb Q$-factorial model
to make the extended divisor classes $\mathbb Q$-Cartier and measure
them by degrees on its fibre curves.  The variety $Y_e$ is terminal,
so the small
$\mathbb Q$-factorialization theorem, the empty-extraction case of
\cite[Corollary~1.4.3]{BCHM}, supplies a projective small morphism
\[
 g:Y'\longrightarrow Y_e
\]
with $Y'$ $\mathbb Q$-factorial.  Since $g$ is small and $K_{Y_e}$
is $\mathbb Q$-Cartier,
\begin{equation}\label{eq:local-small-crepant}
 K_{Y'}=g^*K_{Y_e};
\end{equation}
in particular $Y'$ is terminal.

The morphism $g$ is an isomorphism over the smooth locus of $Y_e$.
To see this directly, push a relatively ample divisor on $Y'$ down
to a Weil divisor on $Y_e$.  On the smooth locus it is Cartier, and
its pullback equals the original divisor there, since $g$ has no
exceptional divisors.  Positive-dimensional fibres would then have
both positive and zero degree against this divisor, a contradiction.
A proper quasi-finite birational morphism to a normal variety is an
isomorphism.

After shrinking $Y_e$ around $\eta'$, all its singular points lie
on $T'$, and every fibre of $g$ has dimension at most one.  For the
last assertion, a component of $g^{-1}(T')$ with generic fibre
dimension at least two would be a divisor in $Y'$ contracted to a
curve, contradicting smallness; the remaining fibre-dimension jumps
form a closed set which does not contain $\eta'$ and may be removed.

Choose Cartier divisors $\widetilde D_i$ representing the line
bundles $\widetilde L_i$ on the smooth integral variety
$\widetilde Y_e$.  Let $D_i$ be their pushforwards to $Y_e$ and
$D_i'$ their strict transforms on $Y'$.  Each $D_i'$ is
$\mathbb Q$-Cartier.  Let
\[
 C_1,\ldots,C_s
\]
be the distinct irreducible curve components of $g^{-1}(\eta')$,
with their reduced structures, counted over
$\kappa(\eta')=K_T$.

\begin{lemma}\label{local:degree-independence}
The vectors
\[
 \bigl(D_i'\cdot C_1,\ldots,D_i'\cdot C_s\bigr)
       \in\mathbb Q^s,\qquad 1\leq i\leq b,
\]
are linearly independent.  In particular $s\geq b$.
\end{lemma}

\begin{proof}
Suppose that
$D'=\sum_i\lambda_iD_i'$, with $\lambda_i\in\mathbb Q$, has zero
degree on every $C_j$.  We will show that $D'$ is numerically trivial
over a neighbourhood of $\eta'$ and that a multiple therefore
descends to $Y_e$.  A line bundle from $Y_e$ has zero punctured Chern
class, so this descent will contradict the independence of the
original classes unless every $\lambda_i$ is zero.

Choose an algebraic closure of $K_T$.  The geometric irreducible
components of an integral $K_T$-curve $C_j$ are permuted transitively
by the absolute Galois group.  The divisor $D'$ is defined over
$K_T$, so its degrees on these components are equal.  Degree is
preserved under extension of the ground field and adds over the
components of the fundamental cycle.  The equality
$D'\cdot C_j=0$ therefore implies zero degree on every geometric
curve component.  This argument also accounts for any multiplicities
if the whole scheme fibre is nonreduced.

All the reduced geometric components are defined over one finite
extension of $K_T$.  Normalize a dense open of $T'$ in that
extension, and shrink until the resulting curve cover is finite
\'etale.  Take the closures of the geometric
components in the pulled-back projective family.  By generic
flatness and openness of geometric integrality, we may make their
curve families flat with geometrically integral fibres.  Their
degrees against a Cartier multiple of $D'$ are then constant and
equal to zero.  Include any zero-dimensional generic components in
this construction as well; after shrinking, their fibres remain
zero-dimensional.  The union of these closures accounts
set-theoretically for the whole family after another shrinking.
Indeed its complement has constructible image in the base curve,
and this image does not contain the generic point, so its closure
is a finite set.  Thus every curve in every remaining fibre over
$T'$ has degree zero.  There are no vertical curves away from $T'$,
where $g$ is an isomorphism.  Removing the finitely many discarded
points from $Y_e$ proves the desired relative numerical triviality.

Choose a positive integer $m$ such that $L=mD'$ is Cartier.  It is
$g$-nef.  By \eqref{eq:local-small-crepant}, $L-K_{Y'}$ is also
$g$-nef.  It is $g$-big: relative bigness is tested on the generic
fibre, which is a point for this birational morphism.  The relative
base point free theorem for the effective klt pair $(Y',0)$ therefore
makes $L$ semiample over $Y_e$; the saturation hypothesis in its
sub-klt formulation is automatic here
\cite[Theorem~2.1 and Remark~3.16]{FujinoBPF}.

We explain why a multiple descends all the way to $Y_e$.  Its
semiample morphism, followed by Stein factorization, gives
\[
 Y'\xrightarrow{h} W\xrightarrow{k}Y_e,
\]
where $h$ has connected fibres, $W$ is normal, $k$ is projective,
and a positive multiple of $L$ is the pullback of a $k$-ample line
bundle on $W$.  Because $L$ has degree zero on all $g$-vertical
curves, $h$ contracts every curve in every $g$-fibre.  These fibres
have dimension at most one and are connected, the latter because
$g$ is proper birational with normal target.  Their images under
$h$ are consequently single points.  Thus every fibre of $k$ is
zero-dimensional.  The morphism $k$ is finite and birational, and
normality of $Y_e$ makes it an isomorphism.  We have proved that
\begin{equation}\label{eq:local-line-descent}
 \mathcal O_{Y'}(m'D')\simeq g^*M
\end{equation}
for a positive integer $m'$ and a line bundle $M$ on $Y_e$, with
$m'$ also clearing all the $\lambda_i$.

On the smooth open of $Y_e$, both $g$ and the resolution are
isomorphisms.  By construction of the pushforwards $D_i$,
\eqref{eq:local-line-descent} identifies the restriction there of
\[
 \bigotimes_i\widetilde L_i^{\otimes m'\lambda_i}
\]
with the pullback of $M$.  At a general closed point of $T'$, the
bundle $M$ is trivial on a neighbourhood in the base.  Its first
Chern class on the punctured smooth neighbourhood is therefore zero.
Naturality of
$R\pi_{e*}\mathbb Q\to Rj_{e*}\mathbb Q$ says that the image of a
resolution-fibre Chern class is precisely this punctured Chern class
of the same bundle.  Hence
\[
 \sum_i\lambda_i\,r\bigl(c_1(\widetilde L_i)\bigr)=0
\]
along a dense open of $T'$.  The independence established after
Lemma~\ref{local:etale-lines} forces every $\lambda_i$ to be zero.
This proves the lemma.
\end{proof}

\subsection{Restriction of the valuations}

We isolate the residue-field argument that prevents distinct
valuations from merging under restriction.  The \'etale
neighbourhood in this statement is not required to be finite.

\begin{lemma}\label{local:restriction-injective}
Let $q:Y_e\to Y$ be a dominant \'etale morphism of
normal integral complex varieties, let $\eta\in Y$, and fix
$\eta'\in Y_e$ above $\eta$ such that
$\kappa(\eta)\to\kappa(\eta')$ is an isomorphism.  Suppose that $Y$
is $\mathbb Q$-Gorenstein.  On normalized divisorial valuations of
$\mathbb C(Y_e)$ with centre the closure of $\eta'$, restriction to
$\mathbb C(Y)$ is already normalized, remains divisorial, and is
injective.  Moreover
\[
 a(w;Y_e,0)=a(w|_{\mathbb C(Y)};Y,0).
\]
\end{lemma}

\begin{proof}
The function-field extension
$\mathbb C(Y_e)/\mathbb C(Y)$ is finite, because the dominant
\'etale morphism is quasi-finite.  The elementary
finite-extension inequalities for valuations imply that the
restriction of a nontrivial discrete valuation is nontrivial and
discrete, that its value group has finite index in the original
group, and that the residue-field extension is finite
\cite[Tag~0ASH]{Stacks}.  Consequently, if $w$ is divisorial, the
normalized restriction $v$ has residue transcendence degree
$\dim Y-1$ and is divisorial.  One can see the last assertion without
any choice of a completion: choose $\dim Y-1$ value-zero rational
functions with algebraically independent residues, resolve their
maps to $\mathbb P^1$ on a proper birational model of $Y$, and then
resolve that model.  The centre of $v$ has dimension at least
$\dim Y-1$, by those residues, and is not the generic point, since
$v$ is nontrivial.  Thus it is a divisor.

Choose a projective birational morphism $\mu:Z\to Y$ with $Z$ smooth
and carrying the prime divisor $E$ with $v=\operatorname{ord}_E$,
and form
\[
 \begin{tikzcd}[column sep=large,row sep=large]
 Z_e=Z\times_Y Y_e \arrow[r,"p"] \arrow[d,"\mu_e"']
     & Z \arrow[d,"\mu"] \\
 Y_e \arrow[r,"q"'] & Y.
 \end{tikzcd}
\]
The scheme $Z_e$ is smooth, because $p$ is
\'etale.  It is also integral with function field
$\mathbb C(Y_e)$.  Indeed, its irreducible components are disjoint
and open, and each has nonempty open image in $Z$.  Every component
therefore meets the inverse image of a dense open $U\subset Y$
over which $\mu$ is an isomorphism.  That inverse image is the
nonempty integral open $q^{-1}(U)\subset Y_e$.  There can be only
one component.  The projective map $\mu_e$ is an isomorphism over
$q^{-1}(U)$, proving the function-field assertion.

By properness of $\mu_e$, the valuation $w$ has a centre $x$ on
$Z_e$.  Its image on $Z$ is $\eta_E$, since its restriction is a
positive multiple of $v$, and its image on $Y_e$ is $\eta'$.
The points lying above this specified pair are the points of
\begin{equation}\label{eq:local-residue-point}
 \operatorname{Spec}\bigl(
   \kappa(\eta_E)\otimes_{\kappa(\eta)}\kappa(\eta')\bigr)
       =\operatorname{Spec}\kappa(\eta_E).
\end{equation}
There is therefore exactly one possible point $x$.  Since $p$ is
\'etale, its local ring
$\mathcal O_{Z_e,x}$ is a discrete valuation ring, and its maximal
ideal is generated by a uniformizer of
$\mathcal O_{Z,\eta_E}$.  Thus the ramification index is one.
Any normalized valuation of $\mathbb C(Y_e)$ centred at this DVR
equals its order: every nonzero element of the fraction field is a
power of a uniformizer times a unit of the DVR, and a centred
valuation is zero on such units.  It follows that $w$ is uniquely
determined by $v$ and $\eta'$.  Its restriction is exactly $v$,
not a larger integral multiple.  This proves both injectivity and
the normalization assertion.

Finally choose compatible canonical divisors in the displayed
\'etale square.  Then
\[
 K_{Z_e}-\mu_e^*K_{Y_e}
       =p^*(K_Z-\mu^*K_Y).
\]
The coefficient at the divisor with generic point $x$ is the
coefficient at $E$, because its ramification index is one.  Adding
one to these coefficients proves equality of log discrepancies.
\end{proof}

\begin{proof}[Proof of Proposition~\ref{prop:local-valuations}]
If $b_T=0$ there is nothing to prove.  Otherwise apply
Lemmas~\ref{local:generic-lines} and~\ref{local:etale-lines} and use
the small model $g:Y'\to Y_e$ constructed above.
Lemma~\ref{local:degree-independence} gives at least $b_T$ distinct
arithmetic curve components $C_i$ of $g^{-1}(\eta')$.  Their closures
$S_i\subset Y'$ are distinct surfaces.  They have codimension two,
so terminality makes $Y'$ smooth at their generic points.  The
ordinary blowup of each smooth generic centre $S_i$ has a unique
exceptional prime there.  Let $w_i$ be its normalized divisorial
valuation.  The codimension-two blowup formula gives
\[
 a(w_i;Y',0)=2.
\]
Their centres $S_i$ distinguish them.  By
\eqref{eq:local-small-crepant}, they have the same log discrepancy
over $Y_e$, and their centres on $Y_e$ are exactly $T'$.

Lemma~\ref{local:restriction-injective} now gives distinct normalized
divisorial valuations
$v_i=w_i|_{\mathbb C(Y)}$, each with centre $T$ and with
\[
 a(v_i;Y,0)=2.
\]
Choose a prime component $D$ of $\Theta$ with positive coefficient
and containing $T$.  Since $Y$ is $\mathbb Q$-factorial, a positive
multiple of $D$ is Cartier near $\eta_T$, with a local equation in
the maximal ideal of $\mathcal O_{Y,\eta_T}$.  A valuation centred
there has strictly positive value on this equation.  Effectivity of
$\Theta$ consequently gives $v_i(\Theta)>0$, and hence
\[
 a(v_i;Y,\Theta)=2-v_i(\Theta)<2.
\]
The centres distinguish the valuations arising from different
curves.  This completes the proof.
\end{proof}

Summing this proposition over the singular curves contained in $Q$
bounds the final sum in \eqref{eq:branch-before-valuations} by the
codimension-three valuation count.  Lemma~\ref{lem:echoes} makes this
count finite, completing the branch inequality as explained at the end
of Section~\ref{sec:branch-topology}.
